\printconcepts

\exercise{T/F: Every function has an inverse.}{F}

\exercise{In your own words explain what it means for a function to be ``one to one.''}{Answers will vary.}

\exercise{If $(1,10)$ lies on the graph of $y=f(x)$, what can be said about the graph of $y=f^{-1}(x)$?}{The point $(10,1)$ lies on the graph of $y=f^{-1}(x)$ (assuming $f$ is invertible).}

%\exercise{If $(1,10)$ lies on the graph of $y=f(x)$ and $\fp(1) = 5$, what can be said about $y=f^{-1}(x)$?}{The point $(10,1)$ lies on the graph of $y=f^{-1}(x)$ (assuming $f$ is invertible) and $(f^{-1})'(10) = 1/5$.}

\exercise{If a function doesn't have an inverse, what can we do to help it have an inverse?}{Restrict the domain of the original function so that it is one to one.}

\printproblems

\exercisesetinstructions{, given the graph of $f$, sketch the graph of $f^{-1}$.}

\exercise{%
\begin{tikzpicture}[alt={Line segments connecting (-6,1) to (-4,4) to (-2,5) to (0,7).}]
\begin{axis}[width=\linewidth, axis equal,axis x line=middle, axis y line=middle,xmin=-9, xmax=9, ymin=-8.25, ymax=8.25,name=myplot,yticklabel style={{font=\scriptsize}},xticklabel style={{font=\scriptsize}}, xtick={-9,-8,-7,-6,-5,-4,...,5,6,7,8,9}, ytick={-8,-7,-6,-5,-4,...,5,6,7,8},%xticklabel={-8,-7,-6,-5,-4,...,5,6,7,8},
,ymajorgrids=true, xmajorgrids=true, yminorgrids=true, xminorgrids=true,major grid style={line width=.2pt,draw=gray!90}]
\draw[draw={\colorone},thick] (axis cs:-6,1)--(axis cs:-4,4)--(axis cs:-2,5)--(axis cs:0,7);
\addplot[draw={\colorone},fill={\colorone},mark=*,only marks] coordinates {(-6,1)(-4,4)(-2,5)(0,7)};
\node[above] at (axis cs:-4,5) {$f(x)$};
\end{axis}
\node [right] at (myplot.right of origin) {\scriptsize $x$};
\node [above] at (myplot.above origin) {\scriptsize $y$};
\end{tikzpicture}%
}{%
\begin{tikzpicture}[alt={Line segments connecting (-6,1) to (-4,4) to (-2,5) to (0,7).  A different set of line segments connect (1,-6), to (4,-4) to (5,-2) to (7,0).}]
\begin{axis}[width=.8\linewidth, axis equal,axis x line=middle, axis y line=middle,xmin=-9, xmax=9, ymin=-8.25, ymax=8.25,name=myplot,yticklabel style={{font=\scriptsize}},xticklabel style={{font=\scriptsize}}, xtick={-9,-8,-7,-6,-5,-4,...,5,6,7,8,9}, ytick={-8,-7,-6,-5,-4,...,5,6,7,8},%xticklabel={-8,-7,-6,-5,-4,...,5,6,7,8},
,ymajorgrids=true, xmajorgrids=true, yminorgrids=true, xminorgrids=true,major grid style={line width=.2pt,draw=gray!90}]
\draw[draw={\colorone},thick] (axis cs:-6,1)--(axis cs:-4,4)--(axis cs:-2,5)--(axis cs:0,7);
\addplot[draw={\colorone},fill={\colorone},mark=*,only marks] coordinates {(-6,1)(-4,4)(-2,5)(0,7)};
\draw[draw={\colortwo},thick] (axis cs:1,-6)--(axis cs:4,-4)--(axis cs:5,-2)--(axis cs:7,0);
\addplot[mark=*,only marks,draw={\colortwo},fill={\colortwo}] coordinates {(1,-6)(4,-4)(5,-2)(7,0)};
\addplot[domain=-9:9,black, thick,smooth,dashed] {x};
\node[above] at (axis cs:-4,5) {$f(x)$};
\end{axis}
\node [right] at (myplot.right of origin) {\scriptsize $x$};
\node [above] at (myplot.above origin) {\scriptsize $y$};
\end{tikzpicture}%
}

\exercise{%
\begin{tikzpicture}[alt={Line segments connecting (-7,7) to (-5,4) to (-3,3) to (-1,1).}]
\begin{axis}[width=\linewidth, axis equal,axis x line=middle, axis y line=middle,xmin=-9, xmax=9, ymin=-8.25, ymax=8.25,name=myplot,yticklabel style={{font=\scriptsize}},xticklabel style={{font=\scriptsize}}, xtick={-9,-8,...,8,9}, ytick={-8,-7,...,7,8},
yticklabels={-8,-7,...,0,,2,3,...,7,8},
%xticklabel={-8,-7,-6,-5,-4,...,5,6,7,8},
,ymajorgrids=true, xmajorgrids=true, yminorgrids=true, xminorgrids=true,major grid style={line width=.2pt,draw=gray!90}]
\draw[draw={\colorone},thick] (axis cs:-7,7)--(axis cs:-5,4)--(axis cs:-3,3)--(axis cs:-1,1);
\addplot[draw={\colorone},fill={\colorone},mark=*,only marks] coordinates {(-7,7)(-5,4)(-3,3)(-1,1)};
\node[above] at (axis cs:-4,5) {$f(x)$};
\end{axis}
\node [right] at (myplot.right of origin) {\scriptsize $x$};
\node [above] at (myplot.above origin) {\scriptsize $y$};
\end{tikzpicture}%
}{%
\begin{tikzpicture}[alt={Line segments connecting (-7,7) to (-5,4) to (-3,3) to (-1,1).  A different set of line segments connect (7,-7) to (4,-5) to (3,-3) to (1,-1).}]
\begin{axis}[width=.8\linewidth, axis equal,axis x line=middle, axis y line=middle,xmin=-9, xmax=9, ymin=-8.25, ymax=8.25,name=myplot,yticklabel style={{font=\scriptsize}},xticklabel style={{font=\scriptsize}}, xtick={-9,-8,...,8,9}, ytick={-8,-7,...,7,8},
yticklabels={-8,-7,...,0,,2,3,...,7,8},
xticklabels={-9,-8,...,0,,2,3,...,8,9},
,ymajorgrids=true, xmajorgrids=true, yminorgrids=true, xminorgrids=true,major grid style={line width=.2pt,draw=gray!90}]
\draw[draw={\colorone},thick] (axis cs:-7,7)--(axis cs:-5,4)--(axis cs:-3,3)--(axis cs:-1,1);
\addplot[draw={\colorone},fill={\colorone},mark=*,only marks] coordinates {(-7,7)(-5,4)(-3,3)(-1,1)};
\addplot[domain=-9:9,black, thick,smooth,dashed] {x};
\draw[draw={\colortwo},thick] (axis cs:7,-7)--(axis cs:4,-5)--(axis cs:3,-3)--(axis cs:1,-1);
\addplot[draw={\colortwo},fill={\colortwo},mark=*,only marks] coordinates {(7,-7)(4,-5)(3,-3)(1,-1)};
\node[above] at (axis cs:-4,5) {$f(x)$};
\end{axis}
\node [right] at (myplot.right of origin) {\scriptsize $x$};
\node [above] at (myplot.above origin) {\scriptsize $y$};
\end{tikzpicture}%
}

\exercisesetend

\input{exercises/02-07-exset-01}

\exercisesetinstructions{, find a restriction of the domain of the given function on which the function will have an inverse.}

\exercise{$f(x)=\sqrt{16-x^2}$}{$[-4,0]$ or $[0,4]$}

\exercise{$g(x)=\sqrt{x^2-16}$}{$(-\infty,-4]$ or $[4,\infty)$.}

\exercise{$r(t)=t^2-6t+9$}{$(-\infty,3]$ or $[3,\infty)$}

\exercise{$f(x)=\dfrac{1-\sqrt x}{1+\sqrt x}$}{This is one-to-one on its domain of $[0,\infty)$.}

\exercisesetend

\exercisesetinstructions{, find the inverse of the given function.}

\exercise{$f(x)=\dfrac{x+1}{x-2}$}{$f^{-1}(x)=\dfrac{2x+1}{x-1}$}

\exercise{$f(x)=x^2 +4$ for $x\ge 0$}{$f^{-1}(x)=\sqrt{x-4}$}

\exercise{$f(x)=e^{x+3}-2$}{$f^{-1}(x)=\ln(x+2)-3$}

\exercise{$f(x)=\ln(x-5)+1$}{$f^{-1}(x)=e^{x-1}+5$}

\exercisesetend


\exercisesetinstructions{, find the exact value.}

\exercise{$\tan^{-1}(0)$}{$0$}

\exercise{$\tan^{-1}(\tan(\pi/7))$}{$\pi/7$}

\exercise{$\cos(\cos^{-1}(-1/5))$}{$-1/5$}

\exercise{$\sin^{-1}(\sin(8\pi/3))$}{$\pi/3$}

\exercise{$\sin(\tan^{-1}(1))$}{$1/\sqrt2$}

\exercise{$\sec(\sin^{-1}(-3/5))$}{$5/4$}

\exercise{$\cos(\tan^{-1}(3/7))$}{$7/\sqrt{58}$}

\exercise{$\sin^{-1}(-\sqrt3/2)$}{$-\pi/3$}

\exercise{$\cos^{-1}(-\sqrt2/2)$}{$3\pi/4$}

\exercise{$\cos^{-1}(\cos(8\pi/7))$}{$6\pi/7$}

\exercisesetend


\exercisesetinstructions{, simplify the expression.}

\exercise{$\sin\Bigl(\tan^{-1}\frac{x}{\sqrt{4-x^2}}\Bigr)$}{$x/2$}

\exercise{$\tan\Bigl(\sin^{-1}\frac{x}{\sqrt{x^2+4}}\Bigr)$}{$x/2$}

\exercise{$\cos\Bigl(\sin^{-1}\frac{5}{\sqrt{x^2+25}}\Bigr)$}{$\abs x/\sqrt{x^2+25}$}

\exercise{$\cot\Bigl(\cos^{-1}\frac{3}{\sqrt{x}}\Bigr)$}{$3/\sqrt{x-9}$}

\exercisesetend

\exercise{Show that for any $x$ in the domain of $\sec^{-1}$ we have $\sec^{-1}x=\cos^{-1}\frac1x$.}{}

\exercise{Show that for $\abs x\le1$ we have $\cos^{-1}x=\frac\pi2-\sin^{-1}x$.\\
Hint: Recall the cofunction identity $\cos\theta=\sin(\frac\pi2-\theta)$ for all $\theta$.}{}

\exercise{Show that for any $x$ we have $\cot^{-1}x=\frac\pi2-\tan^{-1}x$.}{}

\exercise{Show that for $\abs x\ge1$ we have $\csc^{-1}x=\frac\pi2-\sec^{-1}x$.}{}

\exercise{A mass attached to a spring oscillates vertically about the equilibrium position $y=0$ according to the function $y(t)=e^{-t}(\cos(\sqrt3t)+\frac1{\sqrt3}\sin(\sqrt3t))$. Find the first positive time $t$ for which $y(t)=0$.}{$2\pi/(3\sqrt3)$}


%\exercisesetinstructions{, an invertible function $f(x)$ is given along with a point that lies on its graph. Using \autoref{thm:deriv_inverse_functions}, evaluate $\left(f^{-1}\right)'(x)$ at the indicated value.}

\exercise{$f(x) = 5x+10$\\
Point$=(2,20)$ \\
Evaluate $\left(f^{-1}\right)'(20)$}{$\left(f^{-1}\right)'(20)=\frac1{\fp(2)}=1/5$}

\exercise{$f(x) = x^2-2x+4$, $x\geq 1$\\
Point$=(3,7)$ \\
Evaluate $\left(f^{-1}\right)'(7)$}{$\left(f^{-1}\right)'(7)=\frac1{\fp(3)}=1/4$}

\exercise{$f(x) = \sin 2x$, $-\pi/4\leq x\leq \pi/4$\\
Point$=(\pi/6,\sqrt{3}/2)$ \\
Evaluate $\left(f^{-1}\right)'(\sqrt3/2)$}{$\left(f^{-1}\right)'(\sqrt3/2)=\frac1{\fp(\pi/6)}=1$}

\exercise{$f(x) = x^3-6x^2+15x-2$\\
Point$=(1,8)$ \\
Evaluate $\left(f^{-1}\right)'(8)$}{$\left(f^{-1}\right)'(8)=\frac1{\fp(1)}=1/6$}

\exercise{$\ds f(x) = \frac{1}{1+x^2}$, $x\geq 0$\\
Point$=(1,1/2)$ \\
Evaluate $\left(f^{-1}\right)'(1/2)$}{$\left(f^{-1}\right)'(1/2)=\frac1{\fp(1)}=-2$}

% this is ET
\exercise{$ f(x) = 6e^{3x}$\\
Point$=(0,6)$ \\
Evaluate $\left(f^{-1}\right)'(6)$}{$\left(f^{-1}\right)'(6)=\frac1{\fp(0)}=1/18$}

\exercisesetend

%\exercisesetinstructions{, compute the derivative of the given function.}

\exercise{$h(t)=\sin^{-1}(2t)$}{$h'(t)=\frac2{\sqrt{1-4t^2}}$}

\exercise{$f(t)=\sec^{-1}(2t)$}{$\fp(t)=\frac1{\abs t\sqrt{4t^2-1}}$}

\exercise{$g(x)=\tan^{-1}(2x)$}{$g'(x)=\frac{2}{1+4x^2}$}

\exercise{$f(x)=x\sin^{-1}x$}{$\fp(x)=\frac{x}{\sqrt{1-x^2}}+\sin ^{-1}(x)$}

\exercise{$g(t)=\sin t\cos^{-1}t$}{$g'(t)=\cos^{-1}(t)\cos(t)-\frac{\sin (t)}{\sqrt{1-t^2}}$}

\exercise{$f(t)=e^t\ln t$}{$\fp(t)=\frac{e^t}t+e^t\ln t$}

\exercise{$\ds h(x)=\frac{\sin^{-1}x}{\cos^{-1}x}$}{$h'(x)=\frac{\sin^{-1}x+\cos^{-1}x}{\sqrt{1-x^2}(\cos^{-1}x)^2}$}

\exercise{$g(x)=\tan^{-1}(\sqrt{x})$}{$g'(x)=\frac1{\sqrt{x} (2 x+2)}$}

\exercise{$f(x)=\sec^{-1}(1/x)$}{$\fp(x)=-\frac1{\sqrt{1-x^2}}$}

% todo Exercise 7.2.18 is a pre-duplicate of 7.2.19 and should be changed
\exercise{$f(x)=\sin (\sin^{-1}x)$}{\mbox{}\\[-2\baselineskip]\parbox[t]{\linewidth}{\begin{enumext}
\item	$f(x)=x$, so $\fp(x)=1$
\item	$\fp(x)=\cos(\sin^{-1}x)\frac1{\sqrt{1-x^2}}=1$.
\end{enumext}}}

\exercisesetend

%\exercisesetinstructions{, compute the derivative of the given function in two ways:
\begin{enumerate}
\item By simplifying first, then taking the derivative, and
\item by using the Chain Rule first then simplifying.
\end{enumerate}
Verify that the two answers are the same.}

\exercise{$f(x) =\sin (\sin^{-1}x)$}{\mbox{}\\[-2\baselineskip]\parbox[t]{\linewidth}{\begin{enumext}
\item	$f(x) = x$, so $\fp(x) = 1$
\item	$\fp(x) = \cos(\sin^{-1}x)\frac{1}{\sqrt{1-x^2}} = 1$.
\end{enumext}}}

\exercise{$f(x) =\tan^{-1}(\tan x)$}{\mbox{}\\[-2\baselineskip]\parbox[t]{\linewidth}{\begin{enumext}
\item	$f(x) = x$, so $\fp(x) = 1$
\item	$\fp(x) = \frac{1}{1+\tan^2x}\sec^2 x = 1$
\end{enumext}}}

\exercise{$f(x) =\sin(\cos^{-1} x)$}{\mbox{}\\[-2\baselineskip]\parbox[t]{\linewidth}{\begin{enumext}
\item	$f(x) = \sqrt{1-x^2}$, so $\fp(x) = \frac{-x}{\sqrt{1-x^2}}$
\item	$\fp(x) = \cos (\cos^{-1}x)(\frac{-1}{\sqrt{1-x^2}}) =\frac{-x}{\sqrt{1-x^2}}$
\end{enumext}}}

\exercise{$f(x) =\sin(\tan^{-1} x)$}{\mbox{}\\[-2\baselineskip]\parbox[t]{\linewidth}{\begin{enumext}
\item	$f(x)=\frac x{\sqrt{x^2+1}}$, so $\fp(x)=\frac1{(x^2+1)^{3/2}}$
\item	$\fp(x)=\cos(\tan^{-1}x)(\frac1{x^2+1})=\frac1{\sqrt{x^2+1}}\cdot\frac1{x^2+1}$
\end{enumext}}}

\exercisesetend

%\input{exercises/02-07-exset-05}
%%\input{exercises/02-04-exset-06}
%%\input{exercises/02-04-exset-07}
%\printreview
%\exercise{Find $\dfrac{\dd y}{\dd x}$, where $x^2y-y^2x = 1$. }{ $\frac{\dd y}{\dd x} = \frac{y (y-2 x)}{x (x-2 y)}$}

%\exercise{Find the equation of the line tangent to the graph of $x^2+y^2+xy=7$ at the point $(1,2)$. }{ $y=-\frac45(x-1)+2$}

%\exercise{Let $f(x) = x^3+x$. Evaluate $\ds \lim_{s\to 0} \frac{f(x+s)-f(x)}{s}$.}{ $3x^2+1$}

%%\exercise{Approximate the value of $(3.01)^4$ using the tangent line to $f(x) = x^4$ at $x=3$.}{The tangent line to $f(x) = x^4$ at $x=3$ is $y=108(x-3)+81$; thus $(3.01)^4 \approx y(3.01) = 108(.01)+81 = 82.08$. }
